\chapter{抽象空间中的锥优化与 KKT 理论}

\label{chap:07}

第 \ref{chap:06} 章已经把共轭、双共轭、下确界卷积与 Fenchel--Rockafellar 对偶写成了统一的原--对偶语言，但那一章仍然主要围绕一般复合凸问题展开。进入本章之后，约束的几何角色会被进一步突出：谁被视为“非负对象”，谁就决定了可行方向、乘子的取值空间以及互补松弛的形式。有限维里这类结构常被写成广义不等式、锥约束和 KKT 条件；在无穷维里，它们必须被重新表述为对偶空间、法锥和连续线性算子的命题。

本章按照四个层次推进。第 \ref{sec:7-1} 节先用闭凸锥定义偏序与广义不等式，并证明锥序可以完全由对偶锥刻画，把锥约束同第 \ref{sec:4-3} 节的分离原理连接起来。第 \ref{sec:7-2} 节把抽象锥规划写成统一标准型，给出拉格朗日函数、对偶函数以及与第 \ref{sec:6-4} 节 Fenchel 模板的对应关系，为第 \ref{sec:10-3} 节 ADMM 等算法留出统一接口。第 \ref{sec:7-3} 节讨论广义 Slater 条件，说明严格可行性为何能够消除对偶间隙，并解释当锥的拓扑内部退化时，为什么必须转向第 \ref{sec:4-2} 节中的相对内部或代数内部。最后第 \ref{sec:7-4} 节把这些结果收束成无穷维 KKT 系统，把原可行、对偶可行、驻点条件与互补松弛统一为可检验的最优性证据。

这一章的读法有两个重点。其一，必须始终记住乘子不再默认是 Euclidean 向量；它们往往住在对偶空间里，甚至可能是测度或分布。其二，不能把“强对偶成立”和“KKT 成立”当成无条件口号。和第 \ref{sec:5-1} 节直接法中的极小值存在一样，这里每一个结论都依赖于明确的正则性门槛：闭凸锥、适当的连续性、以及 Slater 型资格条件。第 \ref{chap:08} 章会看到，这些结构在 $\mathbb R^n$ 中几乎全部自动坍缩为 Boyd 读者熟悉的矩阵公式；但这种简化恰恰依赖于本章已经完成的抽象铺垫。

\section{广义不等式与闭凸锥}

\label{sec:7-1}

无穷维锥规划的第一步，不是写算法，也不是直接写 KKT，而是先回答一个更基础的问题：在一般线性空间里，什么才叫“非负”？在欧氏空间里，我们习惯用坐标逐项非负或半正定矩阵来回答这个问题；一旦进入函数空间、测度空间或算子空间，这些坐标化描述就不再可靠。能够跨环境保留下来的，是闭凸锥所诱导的偏序结构。广义不等式的全部语言，正是建立在这个观察之上。

\subsection{锥诱导的偏序}

\begin{definition}[锥诱导偏序与适当锥]\label{def:cone-induced-order}
设 $Y$ 为实 Banach 空间，$K\subset Y$ 为非空闭凸锥，即由定义~\ref{def:convex-cone} 已知满足
\[
\lambda K\subset K,\qquad K+K\subset K,\qquad \lambda\ge 0.
\]
对任意 $u,v\in Y$，定义
\[
u\preceq_K v \quad \Longleftrightarrow \quad v-u\in K.
\]
此时称 $\preceq_K$ 为由 $K$ 诱导的\emph{广义不等式}。若进一步有
\[
K\cap(-K)=\{0\},
\]
则称 $K$ 为\emph{尖锥}；在本章里，闭凸且尖的锥将扮演有限维优化中“正正交体”的抽象替身。
\end{definition}

\begin{definition}[对偶锥]\label{def:dual-cone}
设 $K\subset Y$ 为闭凸锥。定义其\emph{对偶锥}为
\[
K^\ast:=\{y^\ast\in Y^\ast:\langle y^\ast,y\rangle\ge 0,\ \forall y\in K\}.
\]
因此，对偶锥中的元素就是对 $K$ 上所有“非负方向”都取非负值的连续线性泛函。
\end{definition}

\begin{proposition}[对偶锥仍是闭凸锥]\label{prop:dual-cone-closed-convex}
设 $K\subset Y$ 为闭凸锥，则 $K^\ast\subset Y^\ast$ 是凸锥，并且在 $Y^\ast$ 的弱$\ast$拓扑下闭。
\end{proposition}

\begin{proof}
若 $y_1^\ast,y_2^\ast\in K^\ast$ 且 $\alpha,\beta\ge 0$，则对任意 $y\in K$，
\[
\langle \alpha y_1^\ast+\beta y_2^\ast,y\rangle
=
\alpha\langle y_1^\ast,y\rangle+\beta\langle y_2^\ast,y\rangle\ge 0,
\]
故 $\alpha y_1^\ast+\beta y_2^\ast\in K^\ast$，这说明 $K^\ast$ 是凸锥。

再固定任意 $y\in K$，集合
\[
H_y:=\{y^\ast\in Y^\ast:\langle y^\ast,y\rangle\ge 0\}
\]
是弱$\ast$闭半空间，因为映射 $y^\ast\mapsto \langle y^\ast,y\rangle$ 在弱$\ast$拓扑下连续。于是
\[
K^\ast=\bigcap_{y\in K} H_y
\]
是弱$\ast$闭集。
\end{proof}

\begin{proposition}[锥序的对偶刻画]\label{prop:cone-order-dual-characterization}
设 $K\subset Y$ 为闭凸锥，$u,v\in Y$。则下列两条等价：
\[
u\preceq_K v;
\]
\[
\langle y^\ast,u\rangle\le \langle y^\ast,v\rangle,\qquad \forall y^\ast\in K^\ast.
\]
换言之，闭凸锥上的广义不等式完全可以由对偶锥上的正泛函检验。
\end{proposition}

\begin{proof}
若 $u\preceq_K v$，则 $v-u\in K$。任取 $y^\ast\in K^\ast$，由定义~\ref{def:dual-cone}，
\[
\langle y^\ast,v-u\rangle\ge 0,
\]
即
\[
\langle y^\ast,u\rangle\le \langle y^\ast,v\rangle.
\]

反之，设对所有 $y^\ast\in K^\ast$ 都有上述不等式，但 $u\not\preceq_K v$。记
\[
w:=v-u.
\]
则 $w\notin K$。由于 $K$ 是闭凸集，而 $w$ 是其外点，由第 \ref{sec:4-3} 节定理~\ref{thm:hahn-banach-geometric} 的几何分离形式，存在非零泛函 $y^\ast\in Y^\ast$ 与常数 $c$，使
\[
\langle y^\ast,w\rangle<c\le \langle y^\ast,k\rangle,\qquad \forall k\in K.
\]
因为 $0\in K$，故 $c\le 0$。又因为 $tk\in K$ 对任意 $t\ge 0$ 仍成立，上式对每个 $t>0$ 给出
\[
c\le t\langle y^\ast,k\rangle.
\]
令 $t\downarrow 0$ 即得 $\langle y^\ast,k\rangle\ge 0$，故 $y^\ast\in K^\ast$。于是
\[
\langle y^\ast,v-u\rangle=\langle y^\ast,w\rangle<c\le 0,
\]
即
\[
\langle y^\ast,v\rangle<\langle y^\ast,u\rangle,
\]
这与假设矛盾。故必有 $u\preceq_K v$。
\end{proof}

\begin{remark}
命题~\ref{prop:cone-order-dual-characterization} 是本节的核心接口。它告诉我们：锥约束不是另一套和对偶理论平行的语言，而是分离定理在锥环境中的直接应用。后文的锥对偶、互补松弛和广义 KKT，都会把“可行性”重新翻译为“对偶锥中所有正泛函都看不出违约”。
\end{remark}

\subsection{锥不等式与互补关系}

由定义~\ref{def:cone-induced-order}，约束
\[
Ax-b\preceq_K 0
\]
等价于
\[
b-Ax\in K.
\]
因此锥规划中的可行性，本质上是在判断某个\emph{松弛变量}是否属于给定的闭凸锥。与之对应，对偶变量将被要求属于对偶锥 $K^\ast$。这两类对象之间的配对
\[
\langle y^\ast,s\rangle,\qquad y^\ast\in K^\ast,\ s\in K,
\]
天然非负，而当该配对等于零时，就出现了后文 KKT 条件里的互补松弛。

\subsection{典型例子}

\begin{example}[正函数锥与正测度锥]\label{ex:positive-cone-measure-space}
设 $\Omega$ 为紧 Hausdorff 空间，取
\[
Y=C(\Omega),\qquad K:=\{u\in C(\Omega):u(\omega)\ge 0,\ \forall \omega\in \Omega\}.
\]
则 $K$ 是 $C(\Omega)$ 中的闭凸锥。由第 \ref{sec:2-3} 节定理~\ref{thm:riesz-representation}，$Y^\ast$ 可识别为有限符号 Radon 测度空间，而 $K^\ast$ 恰好对应于所有正 Radon 测度。把这个例子放在这里，是为了说明锥乘子完全可以是“正测度”而不是向量；这正是第 \ref{sec:7-4} 节里测度型乘子的原型。
\end{example}

\begin{example}[非负正交体与二阶锥]
在 $Y=\mathbb R^m$ 中，若
\[
K=\mathbb R_+^m,
\]
则 $u\preceq_K v$ 就是坐标逐项不等式；若
\[
K=\{(t,z)\in \mathbb R\times \mathbb R^{m-1}:t\ge \|z\|_2\},
\]
则得到二阶锥不等式。把这个例子放在这里，是为了把 LP 和 SOCP 作为抽象闭凸锥的特例立即回收，并为第 \ref{sec:8-3} 节的有限维坍缩做好准备。
\end{example}

\begin{example}[资源非负约束与影子价格]
在管理科学问题里，$b-Ax\in K$ 常表示“资源剩余是非负的”，其中 $K$ 可以是逐项非负锥、鲁棒安全余量锥，甚至某个函数空间中的正锥。对偶锥中的元素则可解释为资源影子价格或风险罚率。把这个例子放在这里，是为了说明抽象锥语言并不比工程语言更远；它恰好把各种“非负成本”与“边际价格”统一到了同一配对框架中。
\end{example}

\subsection*{有限维对照}

Boyd 在讨论广义不等式时，通常从正正交体、半正定锥或二阶锥直接出发，因此读者很容易把“锥”理解成几类可画图的特殊集合。本节说明，这些例子背后真正稳定的结构是定义~\ref{def:cone-induced-order} 与定义~\ref{def:dual-cone}。有限维里能凭直觉看见的几何，在无穷维里必须依赖命题~\ref{prop:cone-order-dual-characterization} 这样的对偶刻画来替代；但两者说的其实是同一件事。

\subsection*{习题（待补）}

\begin{exercise}[对偶锥检验占位]
补一题：对一个给定的闭凸锥 $K$，显式计算 $K^\ast$，并用命题~\ref{prop:cone-order-dual-characterization} 检验某个广义不等式是否成立。
\end{exercise}

\begin{exercise}[锥约束桥接题占位]
补一题：把一个函数空间中的正锥约束退回到 $\mathbb R^m$ 中的逐项非负约束，比较两种环境下对偶锥与乘子含义的变化。
\end{exercise}


\section{抽象凸优化问题标准型}

\label{sec:7-2}

前一节给出了锥诱导的广义不等式，但若没有统一标准型，后文的 Slater 条件、KKT 系统和第 \ref{sec:10-3} 节 ADMM 都会被迫针对每一种约束重新写一遍。本节的任务就是把它们组织成一个可复用的接口层：原问题如何书写，拉格朗日函数如何定义，对偶函数如何得到，以及这些对象怎样与第 \ref{sec:6-4} 节的 Fenchel--Rockafellar 模板接轨。顺便说明一点：按第 \ref{sec:6-1} 节当前正文的规范术语，这里应当说“Fenchel 共轭”，而不是旧控制卡里遗留的“Fenchel-Moreau 共轭”。

\subsection{抽象锥规划标准型}

\begin{definition}[抽象锥规划标准型]\label{def:abstract-cone-program}
设 $X,Y$ 为实 Banach 空间，$A\colon X\to Y$ 为有界线性算子，$b\in Y$，$K\subset Y$ 为闭凸锥，$f\colon X\to \mathbb R\cup\{+\infty\}$ 为 proper 凸泛函。称问题
\[
\inf_{x\in X} f(x)\qquad \text{subject to}\qquad Ax-b\in -K
\]
为\emph{抽象锥规划标准型}。它的可行集记为
\[
\mathcal F:=\{x\in X:Ax-b\in -K\}=\{x\in X:b-Ax\in K\}.
\]
\end{definition}

\begin{definition}[锥规划的拉格朗日函数]\label{def:cone-lagrangian}
在定义~\ref{def:abstract-cone-program} 的设定下，对任意 $y^\ast\in K^\ast$，定义\emph{拉格朗日函数}
\[
L(x,y^\ast):=f(x)+\langle y^\ast,Ax-b\rangle.
\]
相应的\emph{对偶函数}定义为
\[
q(y^\ast):=\inf_{x\in X}L(x,y^\ast),\qquad y^\ast\in K^\ast.
\]
\end{definition}

\begin{proposition}[对偶函数给出原问题下界]\label{prop:cone-dual-function-lower-bound}
设问题 \eqref{def:abstract-cone-program} 可行，记其最优值为
\[
p^\star:=\inf_{x\in \mathcal F} f(x).
\]
则对任意 $y^\ast\in K^\ast$，都有
\[
q(y^\ast)\le p^\star.
\]
\end{proposition}

\begin{proof}
任取可行点 $x\in \mathcal F$。由可行性，$b-Ax\in K$。又因 $y^\ast\in K^\ast$，故
\[
\langle y^\ast,b-Ax\rangle\ge 0,
\]
亦即
\[
\langle y^\ast,Ax-b\rangle\le 0.
\]
于是
\[
L(x,y^\ast)=f(x)+\langle y^\ast,Ax-b\rangle\le f(x).
\]
对 $x\in\mathcal F$ 取下确界，得
\[
q(y^\ast)=\inf_{x\in X}L(x,y^\ast)\le \inf_{x\in\mathcal F}f(x)=p^\star.
\]
\end{proof}

\subsection{与 Fenchel 模板的对应}

\begin{proposition}[锥标准型到 Fenchel 对偶的接口]\label{prop:abstract-standard-form-to-fenchel}
设 $g\colon Y\to \mathbb R\cup\{+\infty\}$ 为集合 $b-K$ 的指示函数：
\[
g(y):=\iota_{b-K}(y).
\]
则定义~\ref{def:abstract-cone-program} 的原问题可改写为
\[
\inf_{x\in X}\{f(x)+g(Ax)\}.
\]
此外，对任意 $y^\ast\in Y^\ast$，
\[
g^\ast(y^\ast)=\langle y^\ast,b\rangle+\iota_{K^\ast}(y^\ast),
\]
从而第 \ref{sec:6-4} 节定义~\ref{def:fenchel-primal-dual-pair} 的 Fenchel 对偶恰化为
\[
\sup_{y^\ast\in K^\ast}\{-f^\ast(-A^\ast y^\ast)-\langle y^\ast,b\rangle\}.
\]
\end{proposition}

\begin{proof}
由 $g=\iota_{b-K}$ 可知，$g(Ax)=0$ 当且仅当 $Ax\in b-K$，亦即 $Ax-b\in -K$；否则 $g(Ax)=+\infty$。因此
\[
f(x)+g(Ax)
=
\begin{cases}
f(x), & x\in \mathcal F,\\
+\infty, & x\notin \mathcal F.
\end{cases}
\]
故原问题确可写成
\[
\inf_{x\in X}\{f(x)+g(Ax)\}.
\]

接下来计算 $g^\ast$。由共轭定义，
\[
g^\ast(y^\ast)=\sup_{y\in b-K}\langle y^\ast,y\rangle.
\]
任写 $y=b-s$，其中 $s\in K$，便得
\[
g^\ast(y^\ast)=\sup_{s\in K}\{\langle y^\ast,b-s\rangle\}
=
\langle y^\ast,b\rangle+\sup_{s\in K}(-\langle y^\ast,s\rangle).
\]
若 $y^\ast\in K^\ast$，则对任意 $s\in K$ 都有 $-\langle y^\ast,s\rangle\le 0$，而取 $s=0$ 即见上确界等于 $0$；若 $y^\ast\notin K^\ast$，则存在 $s_0\in K$ 使 $\langle y^\ast,s_0\rangle<0$，于是
\[
\sup_{t\ge 0}(-\langle y^\ast,ts_0\rangle)=+\infty.
\]
故
\[
g^\ast(y^\ast)=\langle y^\ast,b\rangle+\iota_{K^\ast}(y^\ast).
\]
最后将其代入第 \ref{sec:6-4} 节命题~\ref{prop:fenchel-weak-duality} 与定理~\ref{thm:fenchel-rockafellar-duality} 所对应的对偶模板，即得所述对偶表达式。
\end{proof}

\begin{remark}\label{remark:standard-form-interface}
命题~\ref{prop:abstract-standard-form-to-fenchel} 解释了本节的真正角色：它不是重复第 \ref{chap:06} 章的对偶理论，而是把锥约束问题准确嵌入 Fenchel 框架。后文第 \ref{sec:7-3} 节的 Slater 条件，本质上是在验证指示函数 $g=\iota_{b-K}$ 的资格条件；第 \ref{sec:7-4} 节的 KKT 系统，则是第 \ref{sec:6-4} 节推论~\ref{cor:fenchel-kkt-subdifferential} 的锥版本翻译。再往后，第 \ref{sec:10-3} 节 ADMM 之所以能统一处理分块变量与一致性约束，也正是因为抽象标准型已经把这些结构预先接口化了。
\end{remark}

\subsection{典型例子}

\begin{example}[积分型目标与算子约束]
设 $X=L^2(0,T)$，$Y=L^2(0,T)$，$A$ 为积分算子或观测算子，$K=L^2_+(0,T)$，并考虑
\[
\inf_{x\in X}\left\{\int_0^T \ell(t,x(t))\,dt:\ Ax-b\in -K\right\}.
\]
把这个例子放在这里，是为了说明抽象标准型不是只为矩阵问题设计的；函数空间中的能量最小化、状态约束和观测一致性，同样可以直接写进定义~\ref{def:abstract-cone-program}。
\end{example}

\begin{example}[Boyd 风格有限维标准型]
当 $X=\mathbb R^n$、$Y=\mathbb R^m$、$K=\mathbb R_+^m$ 时，定义~\ref{def:abstract-cone-program} 退化为
\[
\min f(x)\qquad \text{subject to}\qquad Ax\le b.
\]
把这个例子放在这里，是为了把 Boyd 中最熟悉的标准凸优化模型直接接回本章：抽象标准型并没有改变有限维问题，只是把其真正依赖的锥结构显式写了出来。
\end{example}

\begin{example}[分块变量与一致性约束]
在分布式学习或推荐系统里，常见模型是
\[
\min_{x,z}\ \phi(x)+\psi(z)\qquad \text{subject to}\qquad Bx-Cz=0.
\]
把等式约束写成锥约束时，可取 $Y$ 为相应乘积空间，$K=\{0\}$，并把 $(x,z)$ 视为一个联合变量。把这个例子放在这里，是为了给第 \ref{sec:10-3} 节 ADMM 的原--对偶分裂准备统一输入格式。
\end{example}

\subsection*{有限维对照}

Boyd 书里标准型通常写成“目标函数 + 线性等式/不等式约束”，其后紧跟拉格朗日函数和对偶函数的矩阵公式。本节说明，这种写法在本质上只是定义~\ref{def:abstract-cone-program} 与定义~\ref{def:cone-lagrangian} 的有限维坍缩。矩阵记号看上去更具体，但隐藏了锥、对偶锥和指示函数这三个真正驱动对偶理论的对象。

\subsection*{习题（待补）}

\begin{exercise}[标准型改写占位]
补一题：把一个带不等式约束和等式约束的具体凸优化模型改写为定义~\ref{def:abstract-cone-program} 的标准型，并写出其拉格朗日函数与对偶函数。
\end{exercise}

\begin{exercise}[Fenchel 接口桥接题占位]
补一题：从命题~\ref{prop:abstract-standard-form-to-fenchel} 出发，显式计算一个有限维锥规划的 Fenchel 对偶，并和 Boyd 书中的拉格朗日对偶公式逐项对照。
\end{exercise}


\section{广义 Slater 条件}

\label{sec:7-3}

弱对偶几乎不需要额外条件，但强对偶从来不是免费的。第 \ref{sec:6-4} 节中，Fenchel--Rockafellar 对偶之所以成立，是因为指示函数或代价函数在适当点上具备连续性；对锥规划来说，这种连续性最自然的几何化表达就是严格可行性。Slater 条件的作用，不是给出一个“好看的假设”，而是确保约束边界没有在最优值附近制造病态缺口。

\subsection{严格可行性}

\begin{definition}[广义 Slater 条件]\label{def:generalized-slater-condition}
考虑定义~\ref{def:abstract-cone-program} 的锥规划问题。若存在某个 $x_0\in \operatorname{dom}f$ 使得
\[
b-Ax_0\in \operatorname{int}K,
\]
则称该问题满足\emph{广义 Slater 条件}。也就是说，存在一个可行点，其松弛变量不仅属于锥 $K$，而且落在锥的拓扑内部。
\end{definition}

\begin{proposition}[Slater 条件推出 Fenchel 资格条件]\label{prop:slater-implies-fenchel-qualification}
设 $g=\iota_{b-K}$，其中 $K\subset Y$ 为闭凸锥。若定义~\ref{def:generalized-slater-condition} 成立，则 $g$ 在点 $Ax_0$ 处连续。因而第 \ref{sec:6-4} 节定义~\ref{def:constraint-qualification-fenchel} 的 Fenchel 型资格条件成立。
\end{proposition}

\begin{proof}
由 $b-Ax_0\in \operatorname{int}K$，可知
\[
Ax_0\in b-\operatorname{int}K=\operatorname{int}(b-K).
\]
因此存在 $0$ 的某个邻域 $U\subset Y$，使
\[
Ax_0+U\subset b-K.
\]
于是对任意 $u\in U$，
\[
g(Ax_0+u)=\iota_{b-K}(Ax_0+u)=0=g(Ax_0).
\]
这说明 $g$ 在 $Ax_0$ 的邻域内恒等于常数，因而在该点连续。再结合 $x_0\in\operatorname{dom}f$，便得到第 \ref{sec:6-4} 节定义~\ref{def:constraint-qualification-fenchel} 的资格条件。
\end{proof}

\subsection{相对内部与 core 的替代}

\begin{proposition}[core 与相对内部版本的 Slater 替代]\label{prop:core-slater-variant}
设
\[
C:=b-K.
\]
若存在 $x_0\in \operatorname{dom}f$ 使得
\[
Ax_0\in \operatorname{ri} C,
\]
则 $g=\iota_C$ 在 $\operatorname{aff}C$ 上相对于其仿射拓扑在 $Ax_0$ 处连续。特别地，当环境空间中 $\operatorname{int}K=\varnothing$ 但 $\operatorname{ri}C$ 或 $\operatorname{core}C$ 非空时，可以把资格条件与分离论证限制在 $\operatorname{aff}C$ 内部进行。
\end{proposition}

\begin{proof}
由第 \ref{sec:4-2} 节命题~\ref{prop:ri-convex-properties}，$Ax_0\in \operatorname{ri}C$ 意味着存在 $\operatorname{aff}C$ 中的邻域 $V$ 使
\[
(Ax_0+V)\cap \operatorname{aff}C\subset C.
\]
于是 $g$ 限制在 $\operatorname{aff}C$ 上时，在该邻域内恒等于 $0$，所以相对于仿射拓扑在 $Ax_0$ 处连续。

另一方面，第 \ref{sec:4-2} 节定义~\ref{def:algebraic-interior} 与命题~\ref{prop:core-separation-bridge} 告诉我们，对于凸集而言，core 提供的正是“沿可行仿射方向仍有余量”的几何信息。在有限维环境中，relative interior 与 core 会重合；而在无穷维里，拓扑内部经常退化为空集，此时就只能退回到相对内部或代数内部的语言。
\end{proof}

\begin{theorem}[Slater 条件下的锥规划强对偶]\label{thm:slater-strong-duality-cone}
设 $X,Y$ 为 Banach 空间，$A\colon X\to Y$ 为有界线性算子，$K\subset Y$ 为闭凸锥，$f\colon X\to \mathbb R\cup\{+\infty\}$ 为 proper、凸且下半连续的泛函。考虑问题
\[
p^\star:=\inf_{x\in X}\{f(x):Ax-b\in -K\},
\]
以及对偶问题
\[
d^\star:=\sup_{y^\ast\in K^\ast}\{-f^\ast(-A^\ast y^\ast)-\langle y^\ast,b\rangle\}.
\]
若定义~\ref{def:generalized-slater-condition} 成立，则
\[
p^\star=d^\star.
\]
也就是说，广义 Slater 条件消除了该锥规划的对偶间隙。
\end{theorem}

\begin{proof}
按命题~\ref{prop:abstract-standard-form-to-fenchel}，原问题可写成
\[
\inf_{x\in X}\{f(x)+\iota_{b-K}(Ax)\}.
\]
由命题~\ref{prop:slater-implies-fenchel-qualification}，函数 $g=\iota_{b-K}$ 在某个点 $Ax_0$ 处连续，因此第 \ref{sec:6-4} 节定义~\ref{def:constraint-qualification-fenchel} 的资格条件满足。于是可直接应用第 \ref{sec:6-4} 节定理~\ref{thm:fenchel-rockafellar-duality}，得到
\[
\inf_{x\in X}\{f(x)+\iota_{b-K}(Ax)\}
=
\sup_{y^\ast\in Y^\ast}\{-f^\ast(-A^\ast y^\ast)-\iota_{b-K}^\ast(y^\ast)\}.
\]
再由命题~\ref{prop:abstract-standard-form-to-fenchel} 对 $\iota_{b-K}^\ast$ 的显式计算，
\[
\iota_{b-K}^\ast(y^\ast)=\langle y^\ast,b\rangle+\iota_{K^\ast}(y^\ast),
\]
从而上式右侧化为
\[
\sup_{y^\ast\in K^\ast}\{-f^\ast(-A^\ast y^\ast)-\langle y^\ast,b\rangle\}=d^\star.
\]
因此 $p^\star=d^\star$。
\end{proof}

\begin{remark}\label{remark:interior-vs-core-slater}
本节最容易被误读的一点，是把“严格可行”机械地等同为“拓扑内部非空”。在 $\mathbb R^m$ 或有限维矩阵空间中，这样做通常没有问题；但在无穷维里，很多天然锥的内部都会退化。例如在非原子测度空间上的 $L^p_+$ 正锥通常没有拓扑内部，而在 $C(\Omega)$ 中的正函数锥却有由严格正下界给出的内部点。因此，一旦环境空间改变，就必须重新判断该用 $\operatorname{int}$、$\operatorname{ri}$ 还是 core，而不能把 Boyd 中的有限维 Slater 条件直接照搬。
\end{remark}

\subsection{典型例子}

\begin{example}[函数约束留有严格余量]
设 $X=Y=C(\Omega)$，$K=C(\Omega)_+$，并考虑约束
\[
Ax\le_K b.
\]
若存在 $x_0$ 使得
\[
b(\omega)-Ax_0(\omega)\ge \delta>0,\qquad \forall \omega\in \Omega,
\]
则 $b-Ax_0\in \operatorname{int}K$，故 Slater 条件成立。把这个例子放在这里，是为了给无穷维严格可行性一个真正可见的形状：不是“略大于零”的口头说法，而是统一正下界。
\end{example}

\begin{example}[有限维标准 Slater 条件]
当 $K=\mathbb R_+^m$ 时，定义~\ref{def:generalized-slater-condition} 退化为
\[
Ax_0<b
\]
的逐项严格不等式。把这个例子放在这里，是为了把 Boyd 书中最熟悉的一句话准确嵌入本章抽象语言：有限维 Slater 条件只是一般锥内部条件的一个特殊面貌。
\end{example}

\begin{example}[分布鲁棒约束的缓冲区解释]
在分布鲁棒或风险约束模型里，$b-Ax\in K$ 往往意味着某种安全裕度或概率缓冲区非负。若存在一个设计点 $x_0$ 使缓冲区严格落在锥内部，则对偶变量就不会被边界病态放大。把这个例子放在这里，是为了说明“资格条件为何消除对偶间隙”并非纯理论话题；它对应的是模型仍留有操作余地的实际语义。
\end{example}

\subsection*{有限维对照}

Boyd 在强对偶讨论中常把 Slater 条件写成“存在严格可行点”，因为在有限维里 interior、relative interior 与许多局部正则性经常会自动协调。本节说明，这种便利并不属于抽象锥规划本身，而属于欧氏环境。无穷维里一旦锥的拓扑内部消失，就必须把第 \ref{sec:4-2} 节的相对内部和 core 拉回来，重新判断哪一种“严格可行”才是真正有效的资格条件。

\subsection*{习题（待补）}

\begin{exercise}[Slater 条件占位]
补一题：对一个给定的锥规划，判断定义~\ref{def:generalized-slater-condition} 是否成立，并据此说明能否直接调用定理~\ref{thm:slater-strong-duality-cone}。
\end{exercise}

\begin{exercise}[相对内部桥接题占位]
补一题：构造一个锥在环境空间中内部为空但相对内部非空的例子，比较拓扑内部版本与相对内部版本 Slater 条件的差别。
\end{exercise}


\section{广义 KKT 条件}

\label{sec:7-4}

到这里为止，我们已经有了三块拼图：第 \ref{sec:7-1} 节给出锥与对偶锥，第 \ref{sec:7-2} 节给出抽象标准型，第 \ref{sec:7-3} 节给出强对偶成立的资格条件。KKT 理论所做的，就是把这三块拼图与第 \ref{sec:6-4} 节的次微分最优性条件合成一个统一系统。有限维里这常常被写成若干向量等式；无穷维里真正需要小心的是，乘子所属的空间、法锥的表达方式以及互补松弛究竟对哪一个锥而言成立。

\subsection{KKT 系统的定义}

\begin{definition}[锥规划的 KKT 系统]\label{def:kkt-cone-system}
考虑定义~\ref{def:abstract-cone-program} 的原问题
\[
\inf_{x\in X}\{f(x):Ax-b\in -K\}.
\]
称 $(\bar x,\bar y^\ast)\in X\times Y^\ast$ 满足该问题的\emph{KKT 系统}，若
\[
A\bar x-b\in -K,
\]
\[
\bar y^\ast\in K^\ast,
\]
\[
0\in \partial f(\bar x)+A^\ast \bar y^\ast,
\]
\[
\langle \bar y^\ast,b-A\bar x\rangle=0.
\]
四个条件分别称为原可行性、对偶可行性、驻点条件和互补松弛。
\end{definition}

\begin{proposition}[锥松弛与对偶锥的互补配对]\label{prop:cone-complementarity-dual-pair}
设 $K\subset Y$ 为闭凸锥，$s\in K$，$y^\ast\in Y^\ast$。则下列两条等价：
\[
y^\ast\in N_{-K}(-s);
\]
\[
y^\ast\in K^\ast,\qquad \langle y^\ast,s\rangle=0.
\]
等价地，对任意 $b\in Y$，
\[
y^\ast\in N_{b-K}(b-s)
\quad \Longleftrightarrow \quad
y^\ast\in K^\ast,\ \langle y^\ast,s\rangle=0.
\]
\end{proposition}

\begin{proof}
先证明关于 $-K$ 的表述。若 $y^\ast\in N_{-K}(-s)$，则对任意 $k\in K$，由于 $-k\in -K$，有
\[
\langle y^\ast,(-k)-(-s)\rangle\le 0,
\]
即
\[
\langle y^\ast,s-k\rangle\le 0.
\]
取 $k=0$ 得 $\langle y^\ast,s\rangle\le 0$；取 $k=2s$ 又得 $\langle y^\ast,-s\rangle\le 0$，故
\[
\langle y^\ast,s\rangle=0.
\]
再代回上式，便得
\[
\langle y^\ast,k\rangle\ge 0,\qquad \forall k\in K,
\]
即 $y^\ast\in K^\ast$。

反之，若 $y^\ast\in K^\ast$ 且 $\langle y^\ast,s\rangle=0$，则对任意 $k\in K$，
\[
\langle y^\ast,(-k)-(-s)\rangle
=
\langle y^\ast,s\rangle-\langle y^\ast,k\rangle
=
-\langle y^\ast,k\rangle\le 0,
\]
从而 $y^\ast\in N_{-K}(-s)$。

最后注意到平移不会改变法锥的方向结构：
\[
N_{b-K}(b-s)=N_{-K}(-s).
\]
于是第二个等价关系随即成立。
\end{proof}

\subsection{广义 KKT 定理}

\begin{theorem}[无穷维锥规划的 KKT 定理]\label{thm:infinite-dimensional-kkt}
设 $X,Y$ 为 Banach 空间，$A\colon X\to Y$ 为有界线性算子，$K\subset Y$ 为闭凸锥，$f\colon X\to \mathbb R\cup\{+\infty\}$ 为 proper、凸且下半连续的泛函。假设第 \ref{sec:7-3} 节定义~\ref{def:generalized-slater-condition} 成立，并设 $\bar x\in X$ 与 $\bar y^\ast\in Y^\ast$ 分别为原问题和对偶问题的可行点。则下列两条等价：
\begin{enumerate}
  \item $\bar x$ 与 $\bar y^\ast$ 分别是原问题与对偶问题的最优解；
  \item $(\bar x,\bar y^\ast)$ 满足定义~\ref{def:kkt-cone-system} 的 KKT 条件。
\end{enumerate}
\end{theorem}

\begin{proof}
把原问题写成
\[
\inf_{x\in X}\{f(x)+\iota_{b-K}(Ax)\}.
\]
由第 \ref{sec:7-3} 节定理~\ref{thm:slater-strong-duality-cone}，该问题满足强对偶。于是可应用第 \ref{sec:6-4} 节推论~\ref{cor:fenchel-kkt-subdifferential}：$(\bar x,\bar y^\ast)$ 为原--对偶最优对，当且仅当
\[
-A^\ast \bar y^\ast\in \partial f(\bar x),
\qquad
\bar y^\ast\in \partial \iota_{b-K}(A\bar x).
\]
由第 \ref{sec:5-4} 节推论~\ref{cor:normal-cone-indicator}，
\[
\partial \iota_{b-K}(A\bar x)=N_{b-K}(A\bar x).
\]
再令
\[
s:=b-A\bar x.
\]
则原可行性就是 $s\in K$，而命题~\ref{prop:cone-complementarity-dual-pair} 给出
\[
\bar y^\ast\in N_{b-K}(A\bar x)
\quad \Longleftrightarrow \quad
\bar y^\ast\in K^\ast,\ \langle \bar y^\ast,b-A\bar x\rangle=0.
\]
因此，第 \ref{sec:6-4} 节的 Fenchel 次微分最优性条件恰好等价于
\[
A\bar x-b\in -K,\qquad
\bar y^\ast\in K^\ast,\qquad
0\in \partial f(\bar x)+A^\ast \bar y^\ast,\qquad
\langle \bar y^\ast,b-A\bar x\rangle=0.
\]
这就是定义~\ref{def:kkt-cone-system} 的 KKT 系统，故两条等价。
\end{proof}

\begin{corollary}[KKT 的 Fenchel 次微分写法]\label{cor:kkt-via-fenchel-subdifferential}
在定理~\ref{thm:infinite-dimensional-kkt} 的假设下，KKT 系统等价于
\[
-A^\ast \bar y^\ast\in \partial f(\bar x),
\qquad
\bar y^\ast\in \partial \iota_{b-K}(A\bar x).
\]
因此，锥规划中的 KKT 条件本质上并不是另一套平行理论，而是第 \ref{sec:6-4} 节推论~\ref{cor:fenchel-kkt-subdifferential} 在指标函数 $g=\iota_{b-K}$ 上的具体化。
\end{corollary}

\begin{proof}
这正是定理~\ref{thm:infinite-dimensional-kkt} 证明中的中间等价关系。
\end{proof}

\subsection{典型例子}

\begin{example}[状态约束最优控制中的测度乘子]\label{ex:measure-valued-multiplier}
在连续时间最优控制里，若状态约束写成
\[
x(t)\le c(t),\qquad t\in[0,T],
\]
则对应的乘子往往不再是有限维向量，而是区间 $[0,T]$ 上的正测度。把这个例子放在这里，是为了说明定义~\ref{def:kkt-cone-system} 中的 $\bar y^\ast$ 的确可能住在对偶空间甚至测度空间里；“乘子是测度”不是修辞，而是无穷维 KKT 的标准现象。
\end{example}

\begin{example}[线性规划的向量乘子]
当 $K=\mathbb R_+^m$ 时，定义~\ref{def:kkt-cone-system} 退化为
\[
Ax\le b,\qquad y\ge 0,\qquad 0\in \partial f(x)+A^\top y,\qquad y^\top(b-Ax)=0.
\]
把这个例子放在这里，是为了把 Boyd 书中的 KKT 条件完整回收：有限维里熟悉的乘子向量，只是本节抽象对偶变量的一个最简单特例。
\end{example}

\begin{example}[影子价格解释]
在资源配置模型里，约束 $b-Ax\in K$ 表示容量或安全余量非负，而对偶可行变量 $y^\ast\in K^\ast$ 则可解释为影子价格。互补松弛
\[
\langle y^\ast,b-Ax\rangle=0
\]
意味着“有价格的资源必然被用到边界，有剩余的资源不会产生边际价格”。把这个例子放在这里，是为了给 KKT 系统一个管理科学层面的可解释语义。
\end{example}

\subsection*{有限维对照}

Boyd 第 5 章中的 KKT 理论常写成几行矩阵方程，因此读者很容易把它理解成线性代数技巧。本节说明，这些公式真正依赖的是四个抽象结构：锥可行性、对偶锥、次微分驻点条件和互补松弛。第 \ref{sec:8-1} 节会进一步看到，一旦回到 $\mathbb R^n$，很多拓扑和对偶上的麻烦都会自动消失；但正因为本节已经把一般形式写清楚，第 \ref{sec:10-3} 节算法中的原--对偶分裂才能被视作 KKT 系统的计算实现，而不只是某种经验公式。

\subsection*{习题（待补）}

\begin{exercise}[KKT 系统占位]
补一题：对一个给定的锥规划写出定义~\ref{def:kkt-cone-system} 的四条条件，并判断某个候选原--对偶对是否满足它们。
\end{exercise}

\begin{exercise}[乘子空间桥接题占位]
补一题：比较一个有限维线性规划与一个函数空间约束问题中的乘子空间，说明为什么前者的乘子是向量，而后者可能是测度。
\end{exercise}

